\section{Dielectric response from first principles}\label{app:E.6}

The valence electrons of silicon, eight per primitive cell of volume $\Omega = a^3/4 = \SI{40.05}{\angstrom^3}$, have a plasma frequency $\omega_p = (4\pi\alpha\,\nd_e/m_e)^{1/2} = \SI{16.6}{eV}$, and the loss function $L(\vect{q},\omega)$ of \cref{sec:4.4} has its main peak near that energy at small momentum. Where exactly the peak sits, how sharp it is, and how much spectral weight lies near the gap, where the dark-matter rate is largest for light masses, depend on how the response is computed. \Cref{sec:4.6,sec:4.9} built the hierarchy of many-body approximations: the independent-particle bubble, the random-phase approximation with and without local fields, time-dependent density-functional kernels and the Bethe--Salpeter equation. Each becomes a calculation with its own inputs and cost.

\subsection{Computed quantities}

Every approximation starts from the crystal bubble of \cref{sec:4.6}, the Adler--Wiser form of the independent-particle response~\cite{adler1962quantum,wiser1963dielectric},
\begin{equation}
  R^0_{\vect{G}\vect{G}'}(\vect{q},\omega) = \frac{g_{\mathrm{deg}}}{N\Omega}\sum_{\vect{k}}\sum_{n,m}
  \bigl(f_{n\vect{k}} - f_{m,\vect{k}+\vect{q}}\bigr)\,
  \frac{\rho_{mn}(\vect{k},\vect{q}+\vect{G})^*\,\rho_{mn}(\vect{k},\vect{q}+\vect{G}')}{\omega + \varepsilon_{n\vect{k}} - \varepsilon_{m,\vect{k}+\vect{q}} + \ii\eta_{\mathrm{b}}} ,
  \label{eq:E.6-R0}
\end{equation}
with $\vect{q}$ in the first zone, $N$ the number of $\vect{k}$ points, $g_{\mathrm{deg}} = 2$ for spin, the density vertices $\rho_{mn}$ of \cref{sec:3.2}, and a broadening $\eta_{\mathrm{b}}$ that tends to $0^+$ in the exact expression and is finite on a grid. The random-phase dielectric matrix is $\varepsilon_{\vect{G}\vect{G}'} = \delta_{\vect{G}\vect{G}'} - v_{\mathrm{C}}(\vect{q}+\vect{G})\,R^0_{\vect{G}\vect{G}'}$, the proper response being approximated by the bubble (\cref{sec:4.4}). Time-dependent density-functional theory replaces it by the solution of the matrix Dyson equation
\begin{equation}
  R = R^0 + R^0 f_{\mathrm{Hxc}} R , \qquad (f_{\mathrm{Hxc}})_{\vect{G}\vect{G}'} = v_{\mathrm{C}}(\vect{q}+\vect{G})\,\delta_{\vect{G}\vect{G}'} + (f_{\mathrm{xc}})_{\vect{G}\vect{G}'}(\vect{q}),
  \label{eq:E.6-dyson}
\end{equation}
from which $\varepsilon^{-1} = 1 + v_{\mathrm{C}}R$ in matrix form; with $f_{\mathrm{xc}} = 0$ it reduces to the random-phase result. From the inverse matrix come the macroscopic dielectric function $\varepsilon_{\mathrm{mac}} = 1/(\varepsilon^{-1})_{\vect{0}\vect{0}}$ and the loss function. For a momentum transfer outside the first zone, $\vect{q} = \vect{q}_r + \vect{G}_q$ in the notation of \cref{sec:3.1}, the loss function that enters the rate is the diagonal element of the inverse at $\vect{G}_q$,
\begin{equation}
  L(\vect{q},\omega) = \Imag\bigl[-\bigl(\varepsilon^{-1}\bigr)_{\vect{G}_q\vect{G}_q}(\vect{q}_r,\omega)\bigr] ,
  \label{eq:E.6-L}
\end{equation}
so a calculation at dark-matter momentum transfers of several keV needs the matrix at reciprocal vectors far from the origin, not just its head (\cref{sec:10.2}). The Bethe--Salpeter equation leaves the two-point language altogether.

\subsection{Place on the hierarchy}

On the truncation label $\mathfrak{T}$ of \cref{sec:8.5}, these calculations differ only in the many-body approximation $\mathcal{X}$: IP for \cref{eq:E.6-R0} without screening, RPA for the diagonal of $\varepsilon_{\vect{G}\vect{G}'}$, RPA with local fields for the full matrix, TDDFT for \cref{eq:E.6-dyson} with a kernel, and BSE. All are at $n_B = 0$, since the nuclei are clamped and $\varepsilon_\infty$ is the static-lattice constant (the lattice part of $\varepsilon_0$ in a polar crystal comes from \cref{app:E.4}), and at zeroth thermal order, $n_T = 0$, with zero-temperature occupations on the static lattice; the temperature dependence of the gap enters only through a shift (\cref{app:E.8}). As \cref{sec:4.6,sec:8.5} stress, $\mathcal{X}$ is not an expansion in a small number. Its error is estimated from the spread between approximations, and for that estimate to mean anything, every approximation must be converged with the same numerical parameters.

\subsection{Convergence parameters}

The wavefunction cutoff, the local-field cutoff, the band count, the $\vect{k}$ grid and the gap control \cref{eq:E.6-R0,eq:E.6-dyson}, and they interact.

\emph{The wavefunction cutoff} limits the largest momentum transfer at which $\rho_{mn}$ can be trusted. A density vertex at $\vect{q}+\vect{G}$ is an overlap of two sets of plane-wave coefficients displaced by $\vect{q}+\vect{G}$, and when the displacement approaches the cutoff momentum, $k_{\mathrm{cut}} = (2m_eE_{\mathrm{cut}})^{1/2}$, few coefficients overlap. Within that range the vertex is still only as good as the pseudo-wavefunction is near the nucleus, which is the subject of \cref{app:E.2,app:E.7}.

\emph{The local-field cutoff} sets the size $N_G$ of the matrix $\varepsilon_{\vect{G}\vect{G}'}$, through a separate, smaller energy cutoff on $|\vect{q}+\vect{G}|$. Building the matrix costs $\propto N_G^2$ per transition and inverting it $\propto N_G^3$ per $(\vect{q},\omega)$. The test is the change of the quantity of interest, $\varepsilon_{\mathrm{mac}}$ near the gap or $L$ at the target $\vect{q}$, as $N_G$ grows. At large $|\vect{q}|$ the matrix must contain reciprocal vectors around $\vect{G}_q$, not only around the origin.

\emph{The number of empty bands} truncates the sum over $m$ in \cref{eq:E.6-R0}. Its test is a sum rule: $\int_0^\infty\dd\omega\,\omega\,\varepsilon_2 = (\pi/2)\omega_p^2$ and the same holds for $\omega L$ (\cref{sec:4.4}), and a truncated band sum misses the weight at high energy. A free-electron count gives the scale. Above the bottom of the valence band, $N_{\mathrm{b}}$ bands per spin fill a sphere of radius $k$ with $\Omega k^3/(6\pi^2) = N_{\mathrm{b}}$, so the highest included band lies at $k^2/2m_e$ above the bottom.

\emph{The $\vect{k}$ grid} is usually the bottleneck in a semiconductor~\cite{albrecht1998ab}. The bands disperse strongly, and \cref{eq:E.6-R0} is a sum of sharp poles that must merge into a continuum. The broadening $\eta_{\mathrm{b}}$ must exceed the spacing of the transition energies on the grid, of order $|\nabla_{\vect{k}}(\varepsilon_m - \varepsilon_n)|\,\Delta k$; a spectrum computed with too small an $\eta_{\mathrm{b}}$ for its grid shows spurious peaks. For finite $\vect{q}$, the first-zone part $\vect{q}_r$ must be a difference of grid points, so the $\vect{k}$ grid also fixes the momentum resolution. A coarse grid with a moderate $\eta_{\mathrm{b}}$ suffices where the spectrum is smooth on the scale of $\eta_{\mathrm{b}}$; the region near the gap is not, and it needs much finer sampling, often on a shifted or locally refined grid.

\emph{The gap} enters every denominator. Kohn--Sham gaps are too small (\cref{sec:3.4}), and the usual remedies are a rigid scissor shift $\Delta_{\mathrm{scis}}$ of the empty bands, quasiparticle energies from \cref{app:E.3}, or a functional designed for gaps. A scissor keeps the states and raises the transition energies, and its effect on the $f$-sum check follows from that. The density vertices are unchanged (\cref{app:E.3}), so in the optical limit the spectrum moves rigidly, $\varepsilon_2(\omega) = \varepsilon_2^{\mathrm{KS}}(\omega - \Delta_{\mathrm{scis}})$, and its first moment becomes $\int\dd\omega\,\omega\,\varepsilon_2^{\mathrm{KS}}(\omega - \Delta_{\mathrm{scis}}) = (\pi/2)\omega_p^2 + \Delta_{\mathrm{scis}}\int\dd\omega\,\varepsilon_2^{\mathrm{KS}}$, larger by $\Delta_{\mathrm{scis}}\int\dd\omega\,\varepsilon_2^{\mathrm{KS}}$. The excess is the $f$-sum rule of a nonlocal operator, the scissor, and a band-convergence test on a scissored spectrum subtracts it. Codes that build the optical limit from velocity matrix elements must rescale them as in \cref{eq:E.3-vrescale}; otherwise the spectrum acquires a spurious factor $[(\omega - \Delta_{\mathrm{scis}})/\omega]^2$ and is no longer the rigid shift that unchanged states imply. A density vertex computed directly from plane-wave coefficients at finite $\vect{q}$ needs no such care, since it depends on the states and not on their energies.

\emph{Semicore and core states} matter at large $\omega$ and large $\vect{q}$. The silicon $2p$ and germanium $3d$ shells are frozen into the core by most pseudopotentials. Their contribution to $L$ starts at their binding energies, and their contribution to the $f$-sum adds to the valence one. A loss function for $\omega$ above a few tens of eV, or at $|\vect{q}|$ where core orbitals overlap the probe, needs them in the valence or an all-electron treatment (\cref{sec:10.6,app:E.7}).

\subsection{Local fields and exchange--correlation kernels}

Local fields are the difference between the inverse of the averaged matrix and the average of the inverse (\cref{sec:4.6}). Computationally they cost one matrix inversion per $(\vect{q},\omega)$, and they change the spectrum where the response is inhomogeneous on the scale of the cell. In the optical spectrum of silicon they lower and shift the two main peaks to higher energy, which moves the random-phase result further from experiment~\cite{albrecht1998ab}. In the dark-matter rate of silicon and germanium, by contrast, the difference between the loss function with and without local fields is small~\cite{knapen2021dark}, because the rate integrates over $\vect{q}$ and $\omega$ and is dominated by energies and momenta where the response is smooth.

The adiabatic local-density kernel adds little to the random-phase result for the absorption of semiconductors: its spectra are extremely close to those of the random-phase approximation with local fields~\cite{reining2002excitonic}. For energy-loss spectra of solids the same kernel does well, but largely because the random-phase approximation with local fields already agrees with experiment there~\cite{reining2002excitonic}, which is good news for the loss function of a dark-matter rate. The reason is that the kernel is short-ranged, while the electron--hole attraction that produces excitons is long-ranged. A kernel of the form $f_{\mathrm{xc}}(\vect{r},\vect{r}') = -\alpha_{\mathrm{xc}}/(4\pi|\vect{r} - \vect{r}'|)$, which behaves as $-\alpha_{\mathrm{xc}}/|\vect{q}+\vect{G}|^2$ in reciprocal space, fits the measured optical spectrum of silicon almost as well as the Bethe--Salpeter equation, with quasiparticle energies in the bubble and the empirical value $\alpha_{\mathrm{xc}} = 0.2$~\cite{reining2002excitonic}, a fit for silicon rather than a universal constant. (Ref.~\cite{reining2002excitonic} writes $\alpha$; we add the subscript to keep the fine-structure constant unambiguous.) The cost of any two-point kernel is that of the random-phase calculation: the matrices have the same size, and only the Dyson equation changes. This is the attraction of TDDFT for dark-matter loss functions on fine $(\vect{q},\omega)$ grids, and the reason its kernels are worth testing against the Bethe--Salpeter equation on a few points.

\subsection{Optical conductivity and large momentum}

Bosonic absorption (\cref{sec:10.7}) and the crystal Migdal response (\cref{sec:12.3}) take opposite limits of one calculation. For absorption of a bosonic dark-matter particle the relevant quantity is the optical conductivity $\sigma_1 = \omega\varepsilon_2$ at $\vect{q}\to0$, where the transverse and longitudinal dielectric functions coincide (\cref{sec:4.4}). Computing it needs the head of \cref{eq:E.6-R0} in the optical limit, with the velocity matrix elements and their nonlocal commutator (\cref{app:E.2,app:E.7}), a fine $\vect{k}$ grid near the absorption edge, and the approximation that the energy range demands: near the gap the Bethe--Salpeter equation, well above it the random-phase approximation with local fields. Measured optical data are available for most candidate targets, and where they are, they fix $\sigma_1$ better than any calculation. The calculation is then needed where the data stop: at low temperature, below the measured range, or for anisotropic crystals where only some tensor components were measured.

The opposite limit is the loss function at momenta far beyond the zone, up to tens of keV, and energies up to about \SI{100}{eV}, which the crystal Migdal response of \cref{sec:12.3} convolves with the nuclear motion. Every convergence parameter above is pushed at once. The cutoff must hold momentum components well above the largest $|\vect{q}|$, the band sum must reach \SI{100}{eV} above the gap and beyond, and the semicore shells must be in the valence. In addition, the local-field matrix must be centered on $\vect{G}_q$, which at $|\vect{q}| = \SI{30}{keV}$ lies more than seven shortest reciprocal vectors from the origin. The pseudo-wavefunctions are no longer adequate there, and the all-electron reconstruction of \cref{app:E.2,app:E.7} becomes part of the response calculation. The cost grows with all of these at once, and in practice such a calculation is done on a coarse grid of $|\vect{q}|$ and a few directions, with the convergence tested on the $f$-sum rule at each $\vect{q}$.

\subsection{The Bethe--Salpeter equation}

The Bethe--Salpeter equation of \cref{sec:4.9} is an eigenvalue problem in the space of transitions $(v,c,\vect{k})$. Its Hamiltonian, in the resonant (Tamm--Dancoff) form, is
\begin{equation}
  H^{\mathrm{exc}}_{vc\vect{k},v'c'\vect{k}'} = \bigl(\varepsilon_{c\vect{k}} - \varepsilon_{v\vect{k}}\bigr)\,\delta_{vv'}\delta_{cc'}\delta_{\vect{k}\vect{k}'}
  + g_{\mathrm{deg}}\,K^{\mathrm{x}}_{vc\vect{k},v'c'\vect{k}'} - K^{\mathrm{d}}_{vc\vect{k},v'c'\vect{k}'} ,
  \label{eq:E.6-Hexc}
\end{equation}
with the energies from quasiparticle bands (\cref{app:E.3}), an exchange kernel $K^{\mathrm{x}}$ built from the Coulomb kernel without its $\vect{G} = 0$ component, and a direct kernel $K^{\mathrm{d}}$ built from the statically screened interaction $W^{\mathrm{scr}}$ at the momentum $\vect{k} - \vect{k}'$ of the scattered pair (\cref{sec:4.9}); the factor $g_{\mathrm{deg}}$ holds for singlets. Its eigenvalues are the exciton energies $E_\lambda$ and its eigenvectors the envelopes $A^\lambda_{cv\vect{k}}$, and the macroscopic dielectric function is a sum over $\lambda$ of oscillator strengths $|\sum_{cv\vect{k}}A^\lambda_{cv\vect{k}}\rho_{cv}(\vect{k},\vect{q})^*|^2$ over $\omega - E_\lambda$. At finite $\vect{q}$ the transitions are $(v\vect{k}) \to (c,\vect{k}+\vect{q})$ and the same structure gives the excitonic loss function (\cref{sec:10.5}).

The cost has two parts. Building the kernel requires the screened interaction $\varepsilon^{-1}_{\vect{G}\vect{G}'}(\vect{k}-\vect{k}')$, which is a random-phase calculation of its own, and the matrix elements of $K^{\mathrm{d}}$ between all pairs of transitions, $\propto (N_vN_cN_k)^2$ entries, each a sum over $N_G^2$ reciprocal vectors when the full inverse is used. Solving it costs $\propto (N_vN_cN_k)^3$ for a full diagonalization, or $\propto (N_vN_cN_k)^2$ per step of an iterative (Lanczos--Haydock) method that delivers the spectrum without the eigenvectors. The transition-basis dimension $N_vN_cN_k$ is the number to watch. For the optical spectrum of silicon, Ref.~\cite{albrecht1998ab} used norm-conserving pseudopotentials at \SI{15}{Ry}, quasiparticle corrections interpolated onto the fine grid, and up to 2{,}048 $\vect{k}$ points in the full zone. On an 864-point grid four valence and two conduction bands proved sufficient below \SI{4.5}{eV}, and neglecting the off-diagonal elements of $\varepsilon^{-1}$ in $K^{\mathrm{d}}$ changed nothing visible, so the production run used six bands and a diagonal $\varepsilon^{-1}$~\cite{albrecht1998ab}. The dimension was then $2{,}048\times4\times2 \approx 1.6\times10^4$. A dense diagonalization of that size takes $\sim 10^{13}$ operations, trivial today, and the authors identified the $\vect{k}$ grid, not the cutoff, as the bottleneck~\cite{albrecht1998ab}.

Symmetry can be used to reduce the zone, but the spectrum turned out to be extremely sensitive to inconsistent phases when wavefunctions are rotated, especially for degenerate bands, and only a subgroup of the operations was used safely~\cite{albrecht1998ab}. This is the gauge rule of \cref{app:A.5}. A finite $\vect{k}$ grid overestimates individual structures; the shoulder at \SI{3.8}{eV} in the silicon spectrum came out slightly too strong for that reason~\cite{albrecht1998ab}.

Bound excitons near the gap are harder still, because their envelopes are wide in real space and narrow in $\vect{k}$. The Wannier--Mott relation of \cref{sec:4.9}, $E_x^b = \mu_x\alpha^2/(2\varepsilon_\infty^2)$, with the measured binding energy of \SI{14.7}{meV} and $\varepsilon_\infty = 11.7$, gives $\mu_x \approx 0.15\,m_e$ and a radius $a_x = \varepsilon_\infty a_0 m_e/\mu_x \approx \SI{4.2}{nm}$. Resolving an envelope of that size needs a $\vect{k}$ spacing well below $1/a_x \approx \SI{0.024}{\angstrom^{-1}}$, against a zone of width $|\vect{b}_i| = \sqrt3\,2\pi/a = \SI{2.0}{\angstrom^{-1}}$: close to a hundred points per direction near the band edges, and so of order $10^6$ points if the grid were uniform. Calculations of such excitons therefore use grids refined only around the band extrema, or the effective-mass limit, with the Bethe--Salpeter equation reserved for the continuum and the finite-momentum structure (\cref{sec:10.5}).

\subsection{Molecular \texorpdfstring{$GW$}{GW} and Bethe--Salpeter calculations}

Molecular crystals (\cref{sec:16.4}) are often treated molecule by molecule, with the environment added afterwards. The same $GW$ and Bethe--Salpeter equations are then solved in a basis of atom-centered Gaussian orbitals instead of plane waves. The $GW$ quasiparticle energies of molecules improve the Kohn--Sham gaps to within a few tenths of an eV at a cost that scales as the fourth power of the system size, and the Bethe--Salpeter singlet excitation energies that follow have typical errors of \SIrange{0.1}{0.3}{eV}, with a mean error of \SI{0.2}{eV} on a standard benchmark set of about 200 excitations~\cite{blase2020bethe}. An iterative search for the lowest eigenstates costs the same as in TDDFT, $O(N_{\mathrm{orb}}^4)$ operations with $N_{\mathrm{orb}}$ the number of orbitals, when resolution-of-the-identity techniques are used to build the kernel~\cite{blase2020bethe}. The convergence parameter specific to this setting is the Gaussian basis: diffuse functions are needed for the empty states and for the screening, and results are tested along a family of basis sets of increasing size. The crystal environment enters through its dielectric screening, as a reaction field added to the bare Coulomb interaction~\cite{blase2020bethe}, or through a periodic calculation of the full crystal. \Cref{sec:16.4} compares the periodic result with measurement for anthracene: \SI{3.3}{eV} against the measured \SI{3.1}{eV} for the lowest singlet of the crystal.

\subsection{Silicon}

\begin{workedexample}{The loss function of silicon from Kohn--Sham states}\label{we:E-elf-si}
A published dark-matter calculation of the silicon loss function used a plane-wave basis with a \SI{500}{eV} cutoff, a $\Gamma$-centered $8\times8\times8$ grid of $\vect{k}$ points, 70 bands per spin, a functional designed for band gaps with a scissor shift to the measured gap, and the random-phase dielectric matrix with local fields at all $\vect{q}$ of the grid~\cite{knapen2021dark}. For the rate, the matrix was reduced to an angular average of its diagonal elements over all $\vect{q}_r + \vect{G}$ of equal length~\cite[Eq.~(24)]{knapen2021dark}. We check each setting against the estimates of this section.

The cutoff corresponds to $k_{\mathrm{cut}} = (2m_e\times\SI{500}{eV})^{1/2} = \SI{22.6}{keV}$, the range quoted by the authors~\cite{knapen2021dark}. Kinematics allows momentum transfers up to $2\mchi v_{\mathrm{max}}$, \SI{5.3}{keV} per MeV of dark-matter mass (\cref{sec:1.3}), but the valence density vertices fall off beyond a few times $\alpha m_e = \SI{3.7}{keV}$ (\cref{sec:3.5,sec:10.1}), and the cutoff covers that range, provided the pseudo-wavefunctions are accurate there (\cref{app:E.7}). By the free-electron count above, the 70 bands reach $k = \SI{4.7}{\angstrom^{-1}}$, which is \SI{84}{eV} above the bottom of the valence band, about \SI{70}{eV} above the valence maximum since the occupied bands are about \SI{12}{eV} wide (\cref{sec:3.3}). The authors traced the shortfall of the computed loss function at large energy to this truncation~\cite{knapen2021dark}. An $8^3$ grid has spacing $|\vect{b}_i|/8 = \SI{0.25}{\angstrom^{-1}}$, or \SI{0.49}{keV}, which is the resolution in $\vect{q}_r$; it is adequate for a loss function dominated by energies of several eV, and coarse for the structure within a few tenths of an eV of the gap, where excitonic and phonon-assisted effects live (\cref{sec:13.1}).

The $f$-sum rule is the cleanest single check. With $\omega_p = \SI{16.6}{eV}$ from the valence density, the integral $\int_0^\infty\dd\omega\,\omega L(\vect{q},\omega) = (\pi/2)\omega_p^2$ at small $\vect{q}$ is \SI{433}{eV^2}. A calculation that recovers, say, 0.9 of it is missing a tenth of the weight, all at high energy, and that missing weight is a bound on what the band truncation can do to a rate at large $\omega$.
\end{workedexample}

The optical spectrum gives the complementary test near the gap: only the electron--hole attraction of the Bethe--Salpeter equation restores the measured line shape and absolute intensity~\cite{albrecht1998ab}. The same hierarchy matters for a dark-matter rate in proportion to the weight it gives to energies within an eV or two of the gap, which is why the many-body spread of \cref{sec:4.6,sec:4.9} is quoted as an uncertainty for light masses and becomes negligible for heavier ones.

\subsection{Pitfalls}

The head element $\varepsilon_{\vect{0}\vect{0}}$ used in place of $1/(\varepsilon^{-1})_{\vect{0}\vect{0}}$ drops the local fields, and the element at $\vect{G} = 0$ used in place of $\vect{G}_q$ for a momentum transfer beyond the first zone computes the response at the wrong momentum. A broadening too large for the grid smears the onset, which is the part of the spectrum that light dark matter sees. The response computed here screens the density vertex, the electron--phonon vertex of \cref{app:E.5} and, through $W^{\mathrm{scr}}$, the quasiparticle energies of \cref{app:E.3}, so every one of these errors propagates into the rates built on them.
